% SPDX-License-Identifier: Apache-2.0
\section{A FIFO Between Two Channels}
\label{sec:x-queue}

A channel is a buffer of two. Where a unit's output needs more room
than that, because the unit behind it takes in bursts or holds off
for a while, \code{Fifo} in \code{txhdl\_parts} is the depth: a
channel in, a channel out, and \code{D} words between, \code{D} a
power of two stated with its address width \code{AW}. A word is
taken from the input whenever there is room and the oldest word is
offered on the output whenever there is one, so a word is visible the
cycle after it entered; the pointers are a head and a tail of
\code{AW} bits with a full bit to tell full from empty. The part is
written in the lowered subset, as every part is, and its test in the
crate runs it against a queue through three hundred cycles.

The example runs it four words deep between a source that bursts
nine words and a sink that takes nothing for eight cycles: the four
words fill it, \code{full} rises, the source is held until the sink
starts taking, and every word comes out in order.

\begin{figure}[htbp]
\centering
\input{queue_timing.tex}
\caption{The FIFO: a burst in, the fill and the hold, then a drain;
\code{head}, \code{tail} and \code{full} are its state.}
\label{fig:x-queue}
\end{figure}

\lstinputlisting[caption={\code{ex\_queue.rs}. Run with \code{bazel run //lib/examples:ex\_queue}.},label={lst:x-queue}]{lib/examples/ex_queue.rs}

\lstinputlisting[style=ascii,caption={What \code{ex\_queue} printed when the build ran it: the pointers and the full bit per cycle, then the lowering.},label={lst:x-queue-out}]{out_queue.txt}
